%! AP Statistics — Unit 2 Cover Sheet
\documentclass[10pt]{article}
\usepackage{apstats-article}
\usepackage{apstats-boxes}

%=============================================================================
\begin{document}

% ── Full-bleed navy banner ────────────────────────────────────────────────────
\begin{tikzpicture}[remember picture, overlay]
  \fill[navy] (current page.north west)
    rectangle ([yshift=-1.4in]current page.north east);
\end{tikzpicture}
\vspace{-1.3in}

\begin{center}
  {\color{white}\bfseries\fontsize{28}{34}\selectfont Statistics}\\[8pt]
  {\color{goldacc}\bfseries\Large Shepherd \quad 2026--2027}\\[14pt]
  {\color{white}\bfseries\LARGE Unit 2: Exploring Two-Variable Data}\\[4pt]
\end{center}

\vspace{0.30in}

% ── Unit overview ─────────────────────────────────────────────────────────────
\begin{tcolorbox}[colback=sky,colframe=navy,boxrule=0.9pt,arc=2mm,
    left=4mm,right=4mm,top=3mm,bottom=3mm]
  \textbf{Unit Overview.}\quad
  Unit 2 extends statistical thinking to the relationship between two variables, in
  five lessons. Students begin with two categorical variables, reading two-way tables and
  comparing conditional distributions to decide whether an association exists. They then
  turn to two quantitative variables: reading scatterplots, describing direction, form,
  and strength, and measuring linear strength with the correlation $r$. The unit builds the
  least-squares regression line for prediction, uses residuals to judge its fit, and closes
  by asking when a line is the wrong model --- curved residual plots, outliers, and
  influential points.
\end{tcolorbox}

\vspace{0.22in}

% ── Lessons in this unit ──────────────────────────────────────────────────────
\renewcommand{\arraystretch}{1.6}
\begin{skillbox}[Lessons in This Unit]{goldbox}
\begin{tabularx}{\linewidth}{@{} c >{\bfseries}p{0.42\linewidth} X @{}}
  \textbf{\#} & \textbf{Title} & \textbf{Focus} \\
  \hline
  2.1 & Two Categorical Variables
    & Two-way tables; joint, marginal, and conditional relative frequencies; association. \\
  \hline
  2.2 & Scatterplots and Correlation
    & Explanatory vs.\ response; direction, form, strength; the correlation $r$. \\
  \hline
  2.3 & Regression Lines and Residuals
    & Predicting with $\hat{y}=a+bx$; extrapolation; residual $=y-\hat{y}$. \\
  \hline
  2.4 & Least-Squares Regression
    & $b = r\,s_y/s_x$, $a = \bar{y}-b\bar{x}$; slope and intercept in context; $r^2$ and $s$. \\
  \hline
  2.5 & Departures from Linearity
    & Residual plots; outliers, high-leverage and influential points; transformations. \\
  \hline
\end{tabularx}
\end{skillbox}

\vspace{0.18in}

% ── Standards covered ─────────────────────────────────────────────────────────
\begin{skillbox}[AP Statistics Big Ideas \& Learning Objectives --- Unit 2]{greenbox}
\begin{tabularx}{\linewidth}{@{} >{\bfseries}p{0.14\linewidth} X @{}}
  VAR-1.D & Identify questions about possible relationships in data; apparent
            associations may be random. \\
  UNC-1.P--S & Compare, calculate, and represent statistics for two categorical variables
            (two-way tables, conditional distributions) and bivariate quantitative data. \\
  DAT-1.A--C & Describe a scatterplot; determine and interpret the correlation for a linear
            relationship. \\
  DAT-1.D--F & Predict with a linear regression model; calculate residuals; use residual
            plots to describe form. \\
  DAT-1.G--H & Estimate and interpret the least-squares regression line's slope,
            intercept, and $r^2$. \\
  DAT-1.I--J & Identify influential points; predict using transformed variables. \\
\end{tabularx}
\end{skillbox}

\end{document}
